\section{Evaluation Methodology}
\label{sec:evaluation-methodology}

\subsection{Dataset, Split, and Targets}
\label{subsec:eval-data}

We train on the GN$\Omega$SIS corpus~\cite{ferikoglou2026gnomegasis},
which contains approximately 219K Vitis HLS design points.
The source-resolvable SALT-HLS subset contains 150,026 measured
configurations from 55 kernels. Language-policy training retains two
FPGA targets---ZCU104 (\texttt{xczu7ev-ffvc1156-2-e}) and Alveo U200
(\texttt{xcu200-fsgd2104-2-e})---at 3.33, 5, and 10\,ns.
Their full [BRAM$_{18K}$, DSP, FF, LUT] capacities are
[624, 1728, 460800, 230400] and
[4320, 6840, 2364480, 1182240], respectively.

We freeze one family-disjoint split (seed 123) for all learned
components: 130,310 rows from 50 kernels/18 families for training,
13,011 rows from four kernels/three families for validation, and
6,705 rows from the Serrano Kalman family as the locked test set.
No test or external-evaluation workload contributes to parameter
updates or checkpoint selection. The structural GNN follows the same
family holdout but receives QoR supervision only on XCZU7EV/10\,ns,
deliberately leaving device, clock, and budget conditioning to the
language policy.

For each training $(\text{kernel},\text{device},\text{clock})$ case,
we construct 16 component-wise resource-budget requests with a 5\%
minimum per resource. A request is retained only if at least three
measured designs satisfy all four limits; up to 24 same-context
candidates form the objective pool, from which the best
objective-ranked feasible design supplies supervision. Directive
domains depend only on source metadata and are therefore independent
of which measured designs survive a budget.

Unless stated otherwise, end-to-end QoR is evaluated on ZCU104/10\,ns at full device capacity. We deliberately construct the nine-kernel evaluation suite as a heterogeneous stress set rather than drawing all workloads from a single benchmark family. It combines two GN$\Omega$SIS-origin applications excluded from policy training (Kalman and 3D Rendering) with seven workloads from external sources/suites: GAN, PolyBench Covariance, Jacobi-2D, SYR2K, and TRMM, and CHStone AES and JPEG. This composition introduces both provenance and structural shifts while covering dense neural computation, memory-constrained filtering, regular matrix and stencil
nests, cryptography, codec/control-intensive code, and graphics.
The resulting kernels expose 3-21 directive action sites and markedly different baseline resource pressure and optimization headroom.
Consequently, the evaluation does not assume that kernel size or the number of available transformations predicts optimization difficulty;
instead, it tests whether SALT-HLS can select useful policies across different dependence structures, memory constraints, and opportunities for parallelism.

We separately test zero-shot target transfer on four deliberately
contrasting kernels---compute-parallel GAN, control- and buffer-intensive JPEG, regular affine Covariance, and memory-constrained Kalman.
Starting from ZCU104/10\,ns, we independently change either the clock to the unseen 6.67\,ns period or the device to the unseen Alveo U50 while holding the other target dimension fixed. The 6.67\,ns point ($\approx$150\,MHz) lies between the 10 and 5\,ns clocks observed during training and therefore tests interpolation to an unseen timing target;
U50 tests device transfer at the familiar 10\,ns clock. Neither target participates in training or model selection.


\subsection{Objectives and Metrics}
\label{subsec:eval-objectives}
\label{subsec:eval-metrics}

For resource $r\in\{\mathrm{BRAM},\mathrm{DSP},\mathrm{FF},\mathrm{LUT}\}$,
we define normalized utilization
$u_r(d)=100R_r(d)/C_r$ and the training-time area proxy
%
\[
A(d)=\frac{1}{4}\sum_r u_r(d).
\]
%
A design is feasible iff every measured resource satisfies its
requested component-wise budget. The three policies minimize,
respectively, synthesized latency $L$, area proxy $A$, and
ADP $L\,A$; the implementation ranks ADP as
$\log_2 L+\log_2 A$, which is order-equivalent for positive values.
$A$ is a ranking proxy rather than a claim of physical silicon area,
so resource results always report BRAM, DSP, FF, and LUT separately.

For same-target experiments, speedup is
$L_{\mathrm{base}}/L_{\mathrm{method}}$ relative to the pragma-free
Vitis implementation under the same part and requested clock, and
aggregates use the geometric mean. For zero-shot clock/device transfer,
where a same-target pragma-free baseline is unavailable, we instead use
nominal execution time $T=Lc$ and report
$T_{\mathrm{ZCU104},10\mathrm{ns}}/T_{\mathrm{transfer}}$.
We distinguish malformed/materialization failures, HLS-synthesis
failures, and synthesized but resource-infeasible designs rather than
collapsing them into one failure rate.


\subsection{Training Protocol}
\label{subsec:eval-training}

\paragraph{Stage 1.}
We use the pinned DeepSeek-Coder-6.7B base model~\cite{guo2024deepseekcoder}
with 4-bit NF4 quantization. Each objective has a LoRA
adapter~\cite{hu2022lora} with rank 8, $\alpha=16$, dropout 0.1,
and attention-projection targets. Training uses a maximum sequence
length of 7168, batch size 1 with eight-step gradient accumulation,
two epochs, LoRA learning rate $1.5\times10^{-5}$, special-token
learning rate $5\times10^{-6}$, cosine decay, 5\% warm-up, and seed 123.
Only directive RHS values with non-singleton source domains contribute
to the prediction loss.

\paragraph{Structural encoder and Stage 2.}
The structured-MLIR GNN uses the architecture of
\S\ref{subsubsec:gnn-backbone} with 403-D node and 82-D edge features.
It is trained on standardized absolute latency/resource targets using
Smooth-$L_1$ loss ($\beta=0.5$), with rank and resource auxiliary
weights 0.05 and 0.10, respectively; the learning rate is
$3\times10^{-5}$ for at most 40 epochs. The validation-selected
epoch-9 checkpoint exports the pragma-free multiscale action memory
used by all Stage~2 policies.

Stage~2 initializes from the matching objective-specific Stage~1
adapter and freezes the backbone, LoRA, and special-token embeddings.
Only the four compiler-routed cross-attention modules and gates are
trained---8.68M parameters, approximately 0.13\% of the 6.76B-parameter
model. We use effective batch size 8, cosine decay with 5\% warm-up,
and at most three epochs; objective-specific learning rates and the
final checkpoint are selected exclusively on the frozen validation
split.

\paragraph{Stage 3.}
Stage~3 completes SALT-HLS with reference-based DPO
~\cite{rafailov2023dpo} over whole-design, same-context preference
pairs. Pairs share kernel, device, clock, budget, objective, and
supervised schema and differ at the semantic-action level; their order
is determined by measured objective QoR. We use $\beta=0.5$, a 0.02
chosen-NLL/SFT anchor, effective batch size 8, and 40 optimization
steps. The Stage~1 LoRA remains frozen; cross-attention and its gates
are refined with learning rate $2\times10^{-5}$.


\subsection{Baselines and Synthesis Protocol}
\label{subsec:eval-baselines}

The motivation study evaluates GPT-5.6 Sol, Claude Opus 5, and
Gemini 3.1 Pro Preview in one-pass \emph{free-form},
\emph{action+domain}, and \emph{action+domain+RAG} regimes.
The action-matched variants receive the same source action sites,
finite RHS domains, target constraints, and objective as SALT-HLS;
RAG adds general HLS guidance but no GN$\Omega$SIS QoR,
Pareto configurations, SALT-HLS outputs, or synthesis feedback.

For end-to-end comparison we use AutoDSE
~\cite{sohrabizadeh2020autodse}, HARP-DSE
~\cite{sohrabizadeh2023harp}, LIFT~\cite{prakriya2025lift},
and CollectiveHLS~\cite{ferikoglou2024collectivehls} on common
workloads, and additionally compare with the strongest
action+domain+RAG general-purpose LLM on GAN and Kalman.
We report only method--kernel pairs for which a meaningful common
experiment is available; missing experiments are marked
\emph{n/e} or \emph{TBD} rather than imputed.

Unless explicitly stated otherwise, all reported accelerator QoR is
obtained with Vitis HLS 2021.1 under the requested FPGA part and clock.
The pragma-free implementation synthesized under the same target is
the reference for same-target speedup. Finally, we distinguish
\emph{online HLS feedback}---synthesis whose result affects configuration
selection---from synthesis used only to validate the final emitted design.
SALT-HLS uses no online HLS feedback during directive selection and one final synthesis for QoR measurement.